\documentclass[12pt,a4paper,oneside]{book}

% =====================================================================
% PACOTES ESSENCIAIS
% =====================================================================
\usepackage[portuguese]{babel}
\usepackage[T1]{fontenc}
\usepackage{amsmath}
\usepackage{amssymb}
\usepackage{amsthm}
\usepackage{geometry}
\usepackage{fancyhdr}
\usepackage{graphicx}
\usepackage{xcolor}
\usepackage{listings}
\usepackage{hyperref}
\usepackage{booktabs}
\usepackage{array}
\usepackage{multirow}
\usepackage{makeidx}
\usepackage{tikz}
\usepackage{mathtools}
\usepackage[nottoc,notlof,notlot]{tocbibind}

% =====================================================================
% CONFIGURAÇÕES DE GEOMETRIA
% =====================================================================
\geometry{
    left=2.5cm,
    right=2.5cm,
    top=3cm,
    bottom=3cm,
    headheight=14pt
}

% =====================================================================
% CONFIGURAÇÕES DE CABEÇALHO E RODAPÉ
% =====================================================================
\pagestyle{fancy}
\fancyhf{}
\fancyhead[L]{\textit{\rightmark}}
\fancyhead[R]{\thepage}
\fancyfoot[C]{\small Princípios de Modelagem Matemática}
\renewcommand{\headrulewidth}{0.4pt}
\renewcommand{\footrulewidth}{0.4pt}

% =====================================================================
% METADADOS DO DOCUMENTO
% =====================================================================
\title{\textbf{Princípios de Modelagem Matemática} \\[1cm]
        \Large Fundamentos e Aplicações}
\author{Baseado em Dym, C. L. (2004) \\
        \textit{Principles of Mathematical Modeling}}
\date{\today}

% =====================================================================
% INÍCIO DO DOCUMENTO
% =====================================================================
\begin{document}

% Capa
\maketitle

% =====================================================================
% FRONT MATTER
% =====================================================================
\frontmatter

\chapter*{Prefácio}

Este material apresenta uma introdução ao processo de modelagem matemática,
seguindo a abordagem clássica de Clifford L. Dym em sua obra
\textit{Principles of Mathematical Modeling} (2ª edição, 2004).

O curso é estruturado em capítulos progressivos, começando com os conceitos
fundamentais de análise dimensional e evoluindo para técnicas mais avançadas
de modelagem, validação e interpretação de resultados.

\section*{Estrutura do Curso}

\begin{enumerate}
    \item \textbf{Capítulo 1:} Introdução à Modelagem Matemática
    \item \textbf{Capítulo 2:} Análise Dimensional (\textcolor{red}{$\checkmark$ Completo})
    \item \textbf{Capítulo 3:} Escalas e Similaridade
    \item \textbf{Capítulo 4:} Modelagem de Sistemas Dinâmicos
    \item \textbf{Capítulo 5:} Aplicações Práticas
\end{enumerate}

\chapter*{Como Usar este Material}

Cada capítulo é autossuficiente e contém:

\begin{itemize}
    \item \textbf{Teoria:} Conceitos fundamentais com demonstrações rigorosas
    \item \textbf{Exemplos:} Aplicações práticas com cálculos detalhados
    \item \textbf{Código:} Implementações em Python para experimentação
    \item \textbf{Exercícios:} Problemas propostos com gabarito
\end{itemize}

Para melhor aproveitamento:
\begin{enumerate}
    \item Leia a teoria e acompanhe os exemplos
    \item Execute o código fornecido
    \item Estude as soluções dos exercícios
    \item Resolva os exercícios propostos
    \item Implemente seus próprios problemas
\end{enumerate}

\tableofcontents

% =====================================================================
% MAIN MATTER
% =====================================================================
\mainmatter

% Capítulo 2: Análise Dimensional
% Nota: O arquivo capitulo2.tex é um documento completo independente
% Use 'capitulo2/capitulo2.tex' para compilar o capítulo individualmente
%
% Para uma edição integrada de todos os capítulos, compile cada capítulo
% separadamente e depois combine os PDFs resultantes.

\chapter{Análise Dimensional}

\section*{Nota sobre este capítulo}

Este capítulo (Análise Dimensional) foi compilado como um documento independente.
Para acessar o conteúdo completo, consulte o arquivo \texttt{capitulo2/capitulo2.pdf}.

O arquivo \texttt{capitulo2/capitulo2.tex} contém:

\begin{itemize}
    \item Introdução à análise dimensional
    \item Conceitos fundamentais de dimensões e unidades
    \item Princípio da homogeneidade dimensional
    \item Método básico (algoritmo de 7 passos)
    \item Teorema de Buckingham Pi
    \item Sistemas de unidades (SI e Imperial)
    \item 4 exemplos resolvidos
    \item 7 exercícios com gabarito (3 níveis de dificuldade)
\end{itemize}

\subsection*{Como usar este material}

\begin{enumerate}
    \item Leia o capítulo em \texttt{capitulo2/capitulo2.pdf}
    \item Estude os exemplos práticos em \texttt{capitulo2/examples.py}
    \item Realize os exercícios em \texttt{capitulo2/EXERCICIOS.md}
    \item Consulte \texttt{capitulo2/README.md} para conceitos teóricos adicionais
\end{enumerate}

% =====================================================================
% BACK MATTER - Referências e Apêndices
% =====================================================================
\backmatter

\begin{thebibliography}{99}

\bibitem{Dym2004}
Dym, C. L. (2004).
\textit{Principles of Mathematical Modeling}.
2nd Edition.
Elsevier/Academic Press.
ISBN: 978-0-12-226351-3.

\bibitem{Buckingham1914}
Buckingham, E. (1914).
On Physically Similar Systems.
Physical Review, 4(4), 345-376.
doi: 10.1103/PhysRev.4.345

\bibitem{Taylor1974}
Taylor, E. S. (1974).
\textit{Dimensional Analysis for Engineers}.
2nd Edition.
Clarendon Press, Oxford.

\bibitem{Barenblatt1996}
Barenblatt, G. I. (1996).
\textit{Scaling, Self-similarity, and Intermediate Asymptotics}.
Cambridge University Press.

\bibitem{Szirtes1997}
Szirtes, T., \& Rózsa, P. (1997).
\textit{Applied Dimensional Analysis and Modeling}.
Elsevier.

\end{thebibliography}

\appendix

\chapter{Instalação e Configuração}

\section{Requisitos}

Para executar os exemplos em Python desta obra, você precisará de:

\begin{itemize}
    \item Python 3.8+
    \item NumPy (análise numérica)
    \item Matplotlib (visualização de gráficos)
    \item TeX Live ou MiKTeX (para compilar documentos LaTeX)
\end{itemize}

\section{Instalação Rápida}

\subsection{Python e Dependências}

\begin{verbatim}
pip install numpy matplotlib scipy
\end{verbatim}

\subsection{LaTeX}

\textbf{Linux (Ubuntu/Debian):}
\begin{verbatim}
sudo apt-get install texlive-full
\end{verbatim}

\textbf{macOS:}
\begin{verbatim}
brew install --cask mactex
\end{verbatim}

\textbf{Windows:}
Descarregue e instale MiKTeX de \url{https://miktex.org/}

\section{Estrutura dos Arquivos}

\begin{verbatim}
PrincipioModelagemMatematica/
  main.tex                      (este arquivo)
  capitulo2/
    capitulo2.tex               (conteudo do capitulo)
    dimensional_analysis.py
    examples.py
    README.md
    EXERCICIOS.md
    INDEX.md
  capitulo3/                    (proximos capitulos)
  README.md                     (instrucoes do repositorio)
\end{verbatim}

\chapter{Soluções dos Exercícios}

As soluções completas de todos os exercícios estão disponíveis em
\texttt{capitulo2/EXERCICIOS.md} com gabarito resumido e dicas importantes
para cada nível de dificuldade.

\chapter{Tabela de Símbolos e Grandezas}

\section{Dimensões Fundamentais}

\begin{center}
\begin{tabular}{lll}
\hline
\textbf{Grandeza} & \textbf{Símbolo} & \textbf{Unidade SI} \\
\hline
Comprimento & $L$ & metro (m) \\
Massa & $M$ & quilograma (kg) \\
Tempo & $T$ & segundo (s) \\
Temperatura & $\Theta$ & kelvin (K) \\
Corrente elétrica & $A$ & ampère (A) \\
\hline
\end{tabular}
\end{center}

\section{Dimensões Derivadas Comuns}

\begin{center}
\begin{tabular}{lll}
\hline
\textbf{Quantidade} & \textbf{Símbolo} & \textbf{Dimensão} \\
\hline
Velocidade & $v$ & $L/T$ \\
Aceleração & $a$ & $L/T^2$ \\
Força & $F$ & $ML/T^2$ \\
Energia & $E$ & $ML^2/T^2$ \\
Potência & $P$ & $ML^2/T^3$ \\
Pressão & $p$ & $M/(LT^2)$ \\
Viscosidade dinâmica & $\mu$ & $M/(LT)$ \\
Número de Reynolds & $\text{Re}$ & $1$ (adimensional) \\
\hline
\end{tabular}
\end{center}

\end{document}
